% appendix/toyrcu/toyrcu.tex
% mainfile: ../../perfbook.tex
% SPDX-License-Identifier: CC-BY-SA-3.0

\QuickQuizChapter{chp:app:``Toy'' RCU Implementations}{"玩具"RCU实现}{qqztoyrcu}
%
\Epigraph{男人和男孩的唯一区别是玩具的价格。}
	 {M. H\'ebert}
% https://www.ncbi.nlm.nih.gov/pubmed/11548147

本附录提供了一系列RCU实现，按照解决存在保证问题的复杂程度递增排列。
\Cref{sec:app:toyrcu:Lock-Based RCU} 介绍了一个基于简单锁的基本RCU实现，而
\crefthro{sec:app:toyrcu:Per-Thread Lock-Based RCU}
{sec:app:toyrcu:RCU Based on Quiescent States}
则介绍了一系列基于锁的简单RCU实现、
\IXalt{reference counters}{reference count},
和自由运行计数器。
最后， \cref{sec:app:toyrcu:Summary of Toy RCU Implementations}
提供了总结和期望的RCU属性列表。

本附录中的玩具RCU实现不是为了高性能、实用性或任何类型的生产用途而设计的，\footnote{
	然而，生产级的用户级RCU实现已经可用~\cite{MathieuDesnoyers2009URCU,MathieuDesnoyers2012URCU}.}
而是为了清晰性和教学目的。
尽管如此，即使是这些玩具RCU实现也需要对
\cref{chp:Introduction,%
chp:Hardware and its Habits,%
chp:Tools of the Trade,%
chp:Partitioning and Synchronization Design,%
chp:Deferred Processing}
有透彻的理解才能轻松理解。

\section{基于锁的RCU}
\label{sec:app:toyrcu:Lock-Based RCU}

也许最简单的RCU实现利用了锁，如
\cref{lst:app:toyrcu:Lock-Based RCU Implementation}
(\path{rcu_lock.h} and \path{rcu_lock.c}).

\begin{listing}
\input{CodeSamples/defer/rcu_lock@lock_unlock.fcv}\vspace*{-11pt}\fvset{firstnumber=last}
\input{CodeSamples/defer/rcu_lock@synchronize.fcv}\fvset{firstnumber=auto}
\caption{基于锁的RCU实现}
\label{lst:app:toyrcu:Lock-Based RCU Implementation}
\end{listing}

在这个实现中，\co{rcu_read_lock()}获取一个全局自旋锁，\co{rcu_read_unlock()}释放它，\co{synchronize_rcu()}获取后立即释放它。

因为\co{synchronize_rcu()}在获取（并释放）锁之前不会返回，所以它不能在所有先前的RCU读端临界区完成之前返回，从而忠实地实现了RCU语义。
当然，一次只能有一个RCU读者处于其读端临界区中，这几乎完全违背了RCU的目的\@。
此外，\co{rcu_read_lock()}和\co{rcu_read_unlock()}中的锁操作开销极大，
读端开销从单个\Power{5} CPU上的约100纳秒到64-CPU系统上的超过17\emph{微秒}不等。
更糟糕的是，这些锁操作使得\co{rcu_read_lock()}可能参与\IXpl{deadlock cycle}。
此外，在没有递归锁的情况下，RCU读端临界区不能嵌套，最后，
尽管并发的RCU更新原则上可以由一个共同的\IX{grace period}来满足，
但这个实现将宽限期串行化了，阻止了宽限期共享。

\QuickQuizSeries{%
\QuickQuizB{
	为什么
	\cref{lst:app:toyrcu:Lock-Based RCU Implementation}
	中RCU实现的死锁不会在任何其他RCU实现中也是死锁？
}\QuickQuizAnswerB{
	\begin{fcvref}[ln:app:toyrcu:Deadlock in Lock-Based RCU Implementation]
	假设
	\cref{lst:app:toyrcu:Deadlock in Lock-Based RCU Implementation}
	中的函数\co{foo()}和\co{bar()}从不同的CPU并发调用。
	那么\co{foo()}将在\clnref{foo:acq}获取\co{my_lock()}，
	而\co{bar()}将在\clnref{bar:rrl}获取\co{rcu_gp_lock}。
	\end{fcvref}

\begin{listing}
\begin{fcvlabel}[ln:app:toyrcu:Deadlock in Lock-Based RCU Implementation]
\begin{VerbatimL}[commandchars=\\\[\]]
void foo(void)
{
	spin_lock(&my_lock);		\lnlbl[foo:acq]
	rcu_read_lock();		\lnlbl[foo:rrl]
	do_something();
	rcu_read_unlock();
	do_something_else();
	spin_unlock(&my_lock);
}

void bar(void)
{
	rcu_read_lock();		\lnlbl[bar:rrl]
	spin_lock(&my_lock);		\lnlbl[bar:acq]
	do_some_other_thing();
	spin_unlock(&my_lock);
	do_whatever();
	rcu_read_unlock();
}
\end{VerbatimL}
\end{fcvlabel}
\caption{基于锁的RCU实现中的死锁}
\label{lst:app:toyrcu:Deadlock in Lock-Based RCU Implementation}
\end{listing}

	\begin{fcvref}[ln:app:toyrcu:Deadlock in Lock-Based RCU Implementation]
	当\co{foo()}执行到\clnref{foo:rrl}时，它将尝试获取\co{rcu_gp_lock}，
	而该锁被\co{bar()}持有。
	然后当\co{bar()}执行到\clnref{bar:acq}时，它将尝试获取\co{my_lock}，
	而该锁被\co{foo()}持有。
	\end{fcvref}

	每个函数都在等待对方持有的锁，这就是经典的死锁。

	其他RCU实现在\co{rcu_read_lock()}中既不自旋也不阻塞，因此避免了死锁。
}\QuickQuizEndB
%
\QuickQuizE{
	为什么不简单地在
	\cref{lst:app:toyrcu:Lock-Based RCU Implementation}
	的RCU实现中使用读写锁，以允许RCU读者并行执行？
}\QuickQuizAnswerE{
	事实上确实可以以这种方式使用读写锁。
	然而，教科书上的读写锁存在内存竞争问题，
	因此RCU读端临界区需要相当长才能真正允许并行
	执行~\cite{McKenney03a}。

	另一方面，使用在\co{rcu_read_lock()}中以读方式获取的读写锁
	将避免上面提到的死锁条件。
}\QuickQuizEndE
}

很难想象这个实现在生产环境中有用，尽管它确实具有几乎可以在任何用户级应用程序中实现的优点。此外，每个CPU一个锁或使用读写锁的类似实现已经在2.4 Linux内核中用于生产。

这种每CPU一锁方法的修改版本，但改为每线程一锁，将在下一节中描述。

\section{每线程基于锁的RCU}
\label{sec:app:toyrcu:Per-Thread Lock-Based RCU}

\Cref{lst:app:toyrcu:Per-Thread Lock-Based RCU Implementation}
(\path{rcu_lock_percpu.h} and \path{rcu_lock_percpu.c})
展示了基于每线程一个锁的实现。
\co{rcu_read_lock()}和\co{rcu_read_unlock()}函数分别获取和释放当前线程的锁。
\co{synchronize_rcu()}函数依次获取和释放每个线程的锁。
因此，当\co{synchronize_rcu()}开始时正在运行的所有RCU读端临界区
必须在\co{synchronize_rcu()}返回之前完成。

\begin{listing}
\input{CodeSamples/defer/rcu_lock_percpu@lock_unlock.fcv}\vspace*{-11pt}\fvset{firstnumber=last}
\input{CodeSamples/defer/rcu_lock_percpu@sync.fcv}\fvset{firstnumber=auto}
\caption{每线程基于锁的RCU实现}
\label{lst:app:toyrcu:Per-Thread Lock-Based RCU Implementation}
\end{listing}

这个实现确实具有允许并发RCU读者的优点，也避免了单个全局锁可能出现的死锁条件。
此外，读端开销虽然在大约140纳秒时较高，
但无论CPU数量如何都保持在约140纳秒。
然而，更新端开销从单个\Power{5} CPU上的约600纳秒
到64个CPU上的超过100\emph{微秒}不等。

\QuickQuizSeries{%
\QuickQuizB{
	\begin{fcvref}[ln:defer:rcu_lock_percpu:sync:loop]
	在
	\cref{lst:app:toyrcu:Per-Thread Lock-Based RCU Implementation}
	的\clnrefrange{b}{e}循环中，先获取所有锁，然后再全部释放，
	这样不是更简洁吗？
	毕竟，经过这一更改，将会有一个时间点不存在任何读者，
	从而大大简化了事情。
	\end{fcvref}
}\QuickQuizAnswerB{
	进行这一更改将重新引入死锁，所以不，这样做并不更简洁。
}\QuickQuizEndB
%
\QuickQuizM{
	\cref{lst:app:toyrcu:Per-Thread Lock-Based RCU Implementation}
	中所示的实现是否不会发生死锁？
	为什么会或为什么不会？
}\QuickQuizAnswerM{
	一种死锁情况是在\co{synchronize_rcu()}期间持有一个锁，
	而在RCU读端临界区内又获取同一个锁。
	然而，这种情况在任何正确设计的RCU实现中都可能导致死锁。
	毕竟，\co{synchronize_rcu()}原语必须等待所有
	先前存在的RCU读端临界区完成，
	但如果其中一个临界区正在自旋等待
	执行\co{synchronize_rcu()}的线程所持有的锁，
	我们就有了RCU定义中固有的死锁\@。

	另一种死锁发生在尝试嵌套RCU读端临界区时。
	这种死锁是此实现特有的，可以通过使用递归锁来避免，
	或者使用由\co{rcu_read_lock()}以读方式获取、
	由\co{synchronize_rcu()}以写方式获取的读写锁来避免。

	然而，如果排除上述两种情况，
	这个RCU实现不会引入任何死锁情况。
	这是因为只有在执行\co{synchronize_rcu()}时才会获取其他线程的锁，
	而在这种情况下，锁会立即释放，
	这就阻止了不涉及跨\co{synchronize_rcu()}持有的锁的死锁环，
	而那正是上面的第一种情况。
}\QuickQuizEndM
%
\QuickQuizE{
	\cref{lst:app:toyrcu:Per-Thread Lock-Based RCU Implementation}
	中所示RCU算法的一个优点不就是它只使用广泛可用的原语吗，
	例如POSIX pthreads中的原语？
}\QuickQuizAnswerE{
	这确实是一个优点，但不要忘记仍然需要
	\co{rcu_dereference()}和\co{rcu_assign_pointer()}，
	这意味着\co{rcu_dereference()}需要\co{volatile}操作，
	\co{rcu_assign_pointer()}需要内存屏障。
	当然，许多Alpha CPU对这两个原语都需要内存屏障。
}\QuickQuizEndE
}

这种方法在某些情况下可能有用，因为类似的方法在
Linux 2.4 kernel~\cite{Molnar00a}.

接下来描述的基于计数器的RCU实现克服了基于锁的实现的一些缺点。

\section{简单的基于计数器的RCU}
\label{sec:app:toyrcu:Simple Counter-Based RCU}

一个稍微复杂一些的RCU实现如
\cref{lst:app:toyrcu:RCU Implementation Using Single Global Reference Counter}
(\path{rcu_rcg.h} and \path{rcu_rcg.c}).
\begin{fcvref}[ln:defer:rcu_rcg]
这个实现使用了在\clnref{lock_unlock:grc}定义的全局引用计数器\co{rcu_refcnt}。
\co{rcu_read_lock()}原语原子地递增这个计数器，
然后执行内存屏障以确保RCU读端临界区排在原子递增之后。
类似地，\co{rcu_read_unlock()}执行内存屏障来限定RCU读端临界区，
然后原子地递减计数器。
\co{synchronize_rcu()}原语自旋等待引用计数器达到零，
前后都有\IXpl{memory barrier}。
\clnref{sync:poll}上的\co{poll()}仅提供纯粹的延迟，
从纯RCU语义的角度来看可以省略。
同样，一旦\co{synchronize_rcu()}返回，
所有先前的RCU读端临界区都保证已经完成。
\end{fcvref}

\begin{listing}
\input{CodeSamples/defer/rcu_rcg@lock_unlock.fcv}\vspace*{-11pt}\fvset{firstnumber=last}
\input{CodeSamples/defer/rcu_rcg@sync.fcv}\fvset{firstnumber=auto}
\caption{使用单个全局引用计数器的RCU实现}
\label{lst:app:toyrcu:RCU Implementation Using Single Global Reference Counter}
\end{listing}

与\cref{sec:app:toyrcu:Lock-Based RCU}中所示的基于锁的实现形成愉快对比的是，
这个实现允许RCU读端临界区的并行执行。
与\cref{sec:app:toyrcu:Per-Thread Lock-Based RCU}中所示的
每线程基于锁的实现形成愉快对比的是，
它还允许嵌套。
此外，\co{rcu_read_lock()}原语不可能参与\IXpl{deadlock cycle}，
因为它既不自旋也不阻塞。

\QuickQuiz{
	但如果你在调用\co{synchronize_rcu()}期间持有一个锁，
	然后在RCU读端临界区内获取同一个锁呢？
}\QuickQuizAnswer{
	确实，这会使任何合法的RCU实现都死锁。
	但\co{rcu_read_lock()}\emph{真的}参与了死锁环吗？
	如果你认为它参与了，那么请在查看
	\cref{sec:app:toyrcu:RCU Based on Quiescent States}
	中的RCU实现时再问自己同样的问题。
}\QuickQuizEnd

然而，这个实现仍然有一些严重的缺点。
首先，\co{rcu_read_lock()}和\co{rcu_read_unlock()}中的原子操作仍然相当重，
读端开销从单个\Power{5} CPU上的约100纳秒到64-CPU系统上的
接近40\emph{微秒}不等。
这意味着RCU读端临界区必须非常长，才能获得任何真正的读端并行性。
另一方面，在没有读者的情况下，宽限期在大约40\emph{纳秒}内结束，
比Linux内核中的生产级实现快了好几个数量级。

\QuickQuiz{
	当\co{synchronize_rcu()}包含10毫秒的延迟时，
	宽限期怎么可能在40纳秒内结束？
}\QuickQuizAnswer{
	更新端测试在没有读者的情况下运行，
	所以\co{poll()}系统调用从未被调用。
	此外，实际代码中这个\co{poll()}系统调用被注释掉了，
	以便更好地评估更新端代码的真实开销。
	任何生产用途的代码使用\co{poll()}系统调用会更好，
	但话说回来，生产用途使用本节稍后展示的其他实现会更好。
}\QuickQuizEnd

其次，如果有许多并发的\co{rcu_read_lock()}和\co{rcu_read_unlock()}操作，
\co{rcu_refcnt}上将会有极端的内存竞争，
导致昂贵的缓存未命中。
这前两个缺点在很大程度上违背了RCU的一个主要目的，
即提供低开销的读端同步原语。

最后，大量具有长读端临界区的RCU读者可能会阻止\co{synchronize_rcu()}完成，
因为全局计数器可能永远不会达到零。
这可能导致RCU更新的\IX{starvation}，
这在生产环境中当然是不可接受的。

\QuickQuiz{
	在\cref{lst:app:toyrcu:RCU Implementation Using Single Global Reference Counter}
	的RCU实现中，为什么不简单地让\co{rcu_read_lock()}在并发的
	\co{synchronize_rcu()}等待过长时也等待？
	这样不就能防止\co{synchronize_rcu()}被饿死吗？
}\QuickQuizAnswer{
	虽然这确实会消除饥饿，但也意味着\co{rcu_read_lock()}会自旋或阻塞
	等待写者，而写者又在等待读者。
	如果这些读者中的一个正在尝试获取自旋/阻塞的\co{rcu_read_lock()}
	所持有的锁，我们又会遇到死锁。

	简而言之，这个治疗比疾病更糟糕。
	请参阅\cref{sec:app:toyrcu:Starvation-Free Counter-Based RCU}
	获取正确的解决方案。
}\QuickQuizEnd

因此，仍然很难想象这个实现在生产环境中有用，尽管它比基于锁的机制有更多的潜力，例如作为适用于高压调试环境的RCU实现。下一节描述了一种对写者更友好的引用计数方案的变体。

\section{无饥饿的基于计数器的RCU}
\label{sec:app:toyrcu:Starvation-Free Counter-Based RCU}

\Cref{lst:app:toyrcu:RCU Read-Side Using Global Reference-Count Pair}
(\path{rcu_rcpg.h})
展示了一个RCU实现的读端原语，该实现使用了一对引用计数器（\co{rcu_refcnt[]}），
以及一个从该对中选择一个计数器的全局索引（\co{rcu_idx}）、
一个每线程嵌套计数器（\co{rcu_nesting}）、
一个每线程的全局索引快照（\co{rcu_read_idx}）
和一个全局锁（\co{rcu_gp_lock}），
它们本身展示在
\cref{lst:app:toyrcu:RCU Global Reference-Count Pair Data}中。

\begin{listing}
\input{CodeSamples/defer/rcu_rcpg@define.fcv}
\caption{RCU全局引用计数对数据}
\label{lst:app:toyrcu:RCU Global Reference-Count Pair Data}
\end{listing}

\begin{listing}
\input{CodeSamples/defer/rcu_rcpg@r.fcv}
\caption{使用全局引用计数对的RCU读端}
\label{lst:app:toyrcu:RCU Read-Side Using Global Reference-Count Pair}
\end{listing}

\paragraph{设计}

正是双元素\co{rcu_refcnt[]}数组提供了免于饥饿的能力。
关键点在于\co{synchronize_rcu()}只需要等待先前存在的读者。
如果一个新读者在给定的\co{synchronize_rcu()}实例已经开始执行之后启动，
那么该\co{synchronize_rcu()}实例不需要等待那个新读者。
在任何给定时刻，当一个读者通过\co{rcu_read_lock()}进入其RCU读端临界区时，
它递增\co{rcu_refcnt[]}数组中由\co{rcu_idx}变量指示的元素。
当同一个读者通过\co{rcu_read_unlock()}退出其RCU读端临界区时，
它递减之前递增的那个元素，忽略\co{rcu_idx}值的任何后续变化。

这种安排意味着\co{synchronize_rcu()}可以通过对\co{rcu_idx}的值取补来避免饥饿，
如\co{rcu_idx = !rcu_idx}。
假设\co{rcu_idx}的旧值为零，那么新值为一。
在取补操作之后到达的新读者将递增\co{rcu_refcnt[1]}，
而之前递增了\co{rcu_refcnt[0]}的旧读者在退出其RCU读端临界区时
将递减\co{rcu_refcnt[0]}。
这意味着\co{rcu_refcnt[0]}的值将不再被递增，
因此将单调递减。\footnote{
	这个"单调递减"的说法忽略了一个\IX{race condition}。
	这个竞争条件将由\co{synchronize_rcu()}的代码来处理。
	在此期间，我建议暂时搁置怀疑。}
这意味着\co{synchronize_rcu()}所需做的就是等待\co{rcu_refcnt[0]}的值达到零。

有了这些背景，我们准备好来看实际原语的实现了。

\paragraph{实现}

\co{rcu_read_lock()}原语原子地递增\co{rcu_refcnt[]}对中
由\co{rcu_idx}索引的成员，并在每线程变量\co{rcu_read_idx}中保存该索引的快照。
\co{rcu_read_unlock()}原语然后原子地递减该对中
对应\co{rcu_read_lock()}所递增的那个计数器。
然而，由于每个线程只记住一个\co{rcu_idx}的值，
必须采取额外的措施来允许嵌套。
这些额外措施使用每线程\co{rcu_nesting}变量来跟踪嵌套。

\begin{fcvref}[ln:defer:rcu_rcpg:r:lock]
为了使这一切工作，
\cref{lst:app:toyrcu:RCU Read-Side Using Global Reference-Count Pair}
中\co{rcu_read_lock()}的\clnref{pick}获取当前线程的\co{rcu_nesting}实例，
如果\clnref{if}发现这是最外层的\co{rcu_read_lock()}，
则\clnrefrange{cur:b}{cur:e}获取\co{rcu_idx}的当前值，
将其保存在该线程的\co{rcu_read_idx}实例中，
并原子地递增所选的\co{rcu_refcnt}元素。
无论\co{rcu_nesting}的值是什么，\clnref{inc}都会递增它。
\Clnref{mb}执行内存屏障，以确保RCU读端临界区
不会泄漏到\co{rcu_read_lock()}代码之前。
\end{fcvref}

\begin{fcvref}[ln:defer:rcu_rcpg:r:unlock]
类似地，\co{rcu_read_unlock()}函数在\clnref{mb}处执行内存屏障，
以确保RCU读端临界区不会泄漏到\co{rcu_read_unlock()}代码之后。
\Clnref{nest}获取该线程的\co{rcu_nesting}实例，
如果\clnref{if}发现这是最外层的\co{rcu_read_unlock()}，
则\clnref{idx,atmdec}获取该线程的\co{rcu_read_idx}实例
（由最外层的\co{rcu_read_lock()}保存）并原子地递减
所选的\co{rcu_refcnt}元素。
无论嵌套级别如何，\clnref{decnest}都会递减该线程的\co{rcu_nesting}实例。
\end{fcvref}

\begin{listing}
\input{CodeSamples/defer/rcu_rcpg@sync.fcv}
\caption{使用全局引用计数对的RCU更新}
\label{lst:app:toyrcu:RCU Update Using Global Reference-Count Pair}
\end{listing}

\begin{fcvref}[ln:defer:rcu_rcpg:sync]
\Cref{lst:app:toyrcu:RCU Update Using Global Reference-Count Pair}
(\path{rcu_rcpg.c})
展示了对应的\co{synchronize_rcu()}实现。
\Clnref{acq,rel}获取和释放\co{rcu_gp_lock}，
以防止多于一个并发的\co{synchronize_rcu()}实例。
\Clnref{pick,compl}分别获取\co{rcu_idx}的值并对其取补，
使得后续的\co{rcu_read_lock()}实例将使用\co{rcu_refcnt}的
不同元素而不是之前实例使用的那个。
\Clnrefrange{while:b}{while:e}
然后等待\co{rcu_refcnt}的先前元素达到零，
\clnref{mb2}上的内存屏障确保对\co{rcu_refcnt}的检查
不会被重排序到\co{rcu_idx}取补之前。
\Clnrefrange{mb3}{while2:e}重复这个过程，
\clnref{mb5}确保任何后续的回收操作不会被重排序到
\co{rcu_refcnt}检查之前。
\end{fcvref}

\QuickQuizSeries{%
\QuickQuizB{
	\begin{fcvref}[ln:defer:rcu_rcpg:sync]
	在\cref{lst:app:toyrcu:RCU Update Using Global Reference-Count Pair}
	的\co{synchronize_rcu()}中，
	既然紧接着就有一个自旋锁获取，
	为什么还需要\clnref{mb1}上的内存屏障？
	\end{fcvref}
}\QuickQuizAnswerB{
	自旋锁获取只保证自旋锁的临界区不会"泄漏"到获取之前。
	它绝不保证自旋锁获取之前的代码不会被重排序到临界区中。
	这种重排序可能导致从RCU保护的链表中的移除操作
	被重排序到\co{rcu_idx}取补之后，
	这可能允许新启动的RCU读端临界区看到最近被移除的数据元素。

	留给读者的练习：
	使用Promela/\co{spin}等工具来确定
	\cref{lst:app:toyrcu:RCU Update Using Global Reference-Count Pair}
	中的哪些内存屏障（如果有的话）确实是需要的。
	参见\cref{chp:Formal Verification}了解使用这些工具的信息。
	第一个正确完整的回答将获得致谢。
}\QuickQuizEndB
%
\QuickQuizE{
	为什么在
	\cref{lst:app:toyrcu:RCU Update Using Global Reference-Count Pair}
	中计数器要翻转两次？
	单次翻转-等待循环不够吗？
}\QuickQuizAnswerE{
	\begin{fcvref}[ln:defer:rcu_rcpg]
	两次翻转都是绝对必要的。
	要理解这一点，考虑以下事件序列：
	\begin{sequence}
	\item	\cref{lst:app:toyrcu:RCU Read-Side Using Global Reference-Count Pair}
		中\co{rcu_read_lock()}的\clnref{r:lock:cur:b}获取\co{rcu_idx}，
		发现其值为零。
	\item	\cref{lst:app:toyrcu:RCU Update Using Global Reference-Count Pair}
		中\co{synchronize_rcu()}的\clnref{sync:compl}
		对\co{rcu_idx}的值取补，将其设置为一。
	\item	\co{synchronize_rcu()}的\clnrefrange{sync:while:b}{sync:while:e}
		发现\co{rcu_refcnt[0]}的值为零，因此返回。
		（回忆一下，问题是在问如果省略
		\clnrefrange{sync:mb3}{sync:mb5}会发生什么。）
	\item	\co{rcu_read_lock()}的\clnref{r:lock:set,r:lock:cur:e}
		将值零存储到该线程的\co{rcu_read_idx}实例中，
		并分别递增\co{rcu_refcnt[0]}。
		然后执行进入RCU读端临界区。
		\label{sec:app:toyrcu:rcu_rcgp:RCU Read Side Start}
	\item	另一个\co{synchronize_rcu()}实例再次对\co{rcu_idx}取补，
		这次将其值设置为零。
		因为\co{rcu_refcnt[1]}为零，\co{synchronize_rcu()}立即返回。
		（回忆一下\co{rcu_read_lock()}递增的是
		\co{rcu_refcnt[0]}，而不是\co{rcu_refcnt[1]}！）
		\label{sec:app:toyrcu:rcu_rcgp:RCU Grace Period Start}
	\item	在\cref{sec:app:toyrcu:rcu_rcgp:RCU Grace Period Start}
		中开始的宽限期被允许结束了，
		尽管在\cref{sec:app:toyrcu:rcu_rcgp:RCU Read Side Start}
		中先前开始的RCU读端临界区尚未完成。
		这违反了RCU语义，可能允许更新释放
		RCU读端临界区仍在引用的数据元素。
	\end{sequence}

	留给读者的练习：
	如果\co{rcu_read_lock()}在\clnref{r:lock:cur:b}之后
	被抢占了很长时间（几个小时！\@）会发生什么？
	在这种情况下，这个实现能正确运行吗？
	为什么能或为什么不能？
	第一个正确完整的回答将获得致谢。
	\end{fcvref}
}\QuickQuizEndE
}

这个实现避免了可能出现在所示的单计数器实现中的更新饥饿问题
\cref{lst:app:toyrcu:RCU Implementation Using Single Global Reference Counter}.

\paragraph{讨论}

仍然存在一些严重的缺点。
首先，\co{rcu_read_lock()}和\co{rcu_read_unlock()}中的原子操作仍然相当重。
事实上，它们比
\cref{lst:app:toyrcu:RCU Implementation Using Single Global Reference Counter}
中所示的单计数器变体更复杂，
读端原语在单个\Power{5} CPU上消耗约150纳秒，
在64-CPU系统上几乎40\emph{微秒}。
更新端\co{synchronize_rcu()}原语的开销也更大，
从单个\Power{5} CPU上的约200纳秒到64-CPU系统上的
超过40\emph{微秒}不等。
这意味着RCU读端临界区必须非常长，才能获得任何真正的读端并行性。

其次，如果有许多并发的\co{rcu_read_lock()}和\co{rcu_read_unlock()}操作，
\co{rcu_refcnt}元素上将会有极端的内存竞争，
导致昂贵的缓存未命中。
这进一步延长了提供并行读端访问所需的RCU读端临界区持续时间。
这前两个缺点在大多数情况下违背了RCU的目的。

第三，需要翻转\co{rcu_idx}两次给更新带来了相当大的开销，
特别是当线程数量很多时。

最后，尽管并发的RCU更新原则上可以由一个共同的宽限期来满足，
但这个实现将宽限期串行化了，阻止了宽限期共享。

\QuickQuiz{
	\begin{fcvref}[ln:defer:rcu_rcpg:r]
	既然原子递增和递减开销如此之大，
	为什么不在
	\cref{lst:app:toyrcu:RCU Read-Side Using Global Reference-Count Pair}
	的\clnref{lock:cur:e}上使用非原子递增，
	在\clnref{unlock:atmdec}上使用非原子递减呢？
	\end{fcvref}
}\QuickQuizAnswer{
	使用非原子操作会导致递增和递减丢失，
	进而导致实现失败。
	请参阅\cref{sec:app:toyrcu:Scalable Counter-Based RCU}
	了解在\co{rcu_read_lock()}和\co{rcu_read_unlock()}中
	安全使用非原子操作的方法。
}\QuickQuizEnd

尽管有这些缺点，可以想象这种RCU变体用于小型紧密耦合的多处理器，也许作为一种节省内存的实现，与更复杂的实现保持API兼容性。然而，它不太可能在超过几个CPU的情况下良好扩展。

下一节描述了引用计数方案的另一种变体，提供了大大改进的读端性能和可扩展性。

\section{可扩展的单计数器RCU}
\label{sec:app:toyrcu:Scalable Single-Counter RCU}

\cref{sec:app:toyrcu:Simple Counter-Based RCU}中使用的单一全局计数器
在任何线程执行RCU读端临界区时都将为非零。
假设每个线程花费固定比例的时间在这样的临界区中，
随着线程的增加，全局计数器为非零的概率趋近于一，
导致更新者饥饿。
避免这种饥饿的一种方法是使用一对全局计数器，
如\cref{sec:app:toyrcu:Starvation-Free Counter-Based RCU}中所讨论的，
但这仍然受到内存竞争的影响，如
\cref{fig:count:Atomic Increment Scalability on x86}所示。
另一种方法是使用每线程计数器，
这样无论有多少线程，给定线程的计数器在该线程处于
RCU读端临界区之外时将为零。

\Cref{lst:app:toyrcu:RCU Read-Side Using Per-Thread Reference-Count}
(\path{rcu_rcpl.h})
展示了使用每线程引用计数器的RCU实现的读端原语。\footnote{
	或者，如果你更愿意的话，嵌套级别计数器。}
每线程\co{rcu_nesting[]}计数器的一个好处是
\co{rcu_read_lock()}和\co{rcu_read_unlock()}原语不再执行原子指令，
而是执行每线程计数器的非原子递增和递减。

\begin{listing}
\input{CodeSamples/defer/rcu_rcl@define.fcv}
\caption{RCU每线程引用计数数据}
\label{lst:app:toyrcu:RCU Per-Thread Reference-Count Data}
\end{listing}

\begin{listing}
\input{CodeSamples/defer/rcu_rcl@r.fcv}
\caption{使用每线程引用计数的RCU读端}
\label{lst:app:toyrcu:RCU Read-Side Using Per-Thread Reference-Count}
\end{listing}

\QuickQuiz{
	别开玩笑了！
	我们可以在\co{rcu_read_lock()}中看到\co{atomic_read()}原语！！！
	那你为什么要假装\co{rcu_read_lock()}不包含原子操作？？？
}\QuickQuizAnswer{
	是的，\co{atomic_read()}原语是一个原子操作，
	但它不执行原子机器指令。
	它只是执行从\co{atomic_t}的普通加载。
	它唯一的目的是让编译器的类型检查满意。
	如果Linux内核运行在8位CPU上，它还需要防止"存储撕裂"，
	这可能是由于需要在某些8位系统上用两次八位访问来存储
	一个16位指针而发生的。
	但值得庆幸的是，似乎没有人在8位系统上运行Linux。
}\QuickQuizEnd

\begin{listing}
\input{CodeSamples/defer/rcu_rcl@u.fcv}
\caption{使用每线程引用计数的RCU更新}
\label{lst:app:toyrcu:RCU Update Using Per-Thread Reference-Count}
\end{listing}

\Cref{lst:app:toyrcu:RCU Update Using Per-Thread Reference-Count}
(\path{rcu_rcl.c})
展示了\co{synchronize_rcu()}的实现。
\begin{fcvref}[ln:defer:rcu_rcl:u:sync]
\co{synchronize_rcu()}函数类似于
\cref{lst:app:toyrcu:RCU Implementation Using Single Global Reference Counter}
中所示的，不同之处在于对全局计数器的轮询被
\clnrefrange{loop:b}{loop:e}上依次轮询每个每线程计数器的扫描所替代。
\end{fcvref}

这个RCU实现对其软件环境提出了重要的新要求，即
%
\begin{enumerate*}[(1)]
\item {}必须能够声明每线程变量，
\item {}这些每线程变量必须能从其他线程访问，以及
\item {}必须能够枚举所有线程。
\end{enumerate*}
%
这些要求在几乎所有软件环境中都可以满足，
但通常会导致线程数量有固定的上限。
更复杂的实现可能会避免这种限制，例如使用可扩展的哈希表。
这种实现可能会动态跟踪线程，
例如在它们第一次调用\co{rcu_read_lock()}时添加。

\QuickQuiz{
	好吧，如果我们有$N$个线程，我们可能有$2N$个十毫秒的等待
	（每次\co{flip_counter_and_wait()}调用一组，
	而且这还假设我们对每个线程只等待一次）。
	我们不需要宽限期\emph{更快}地完成吗？
}\QuickQuizAnswer{
	请记住，我们只在给定线程仍处于RCU读端临界区时才等待该线程。
	因此，我们等待$2N$个间隔的唯一方式是
	每个线程恰好在更新迭代到该线程时处于RCU读端临界区中。

	然而，\cref{sec:app:toyrcu:Scalable Counter-Based RCU}
	中的实现做得更好。
}\QuickQuizEnd

这个实现仍然有几个缺点。
首先，\co{synchronize_rcu()}现在必须检查数量随线程数线性增长的变量，
这给具有大量线程的应用程序带来了相当大的开销。

第三，与以前一样，尽管并发的RCU更新原则上可以由一个共同的宽限期来满足，
但这个实现将宽限期串行化了，阻止了宽限期共享。

最后，如文中所述，对每线程变量和枚举线程的需求在某些软件环境中可能有问题。

话虽如此，读端原语的扩展性非常好，执行一组固定的机器指令来访问每线程计数器。

下一节描述了一种算法，可以防止读者使即使最不幸的更新者也饥饿。

\section{可扩展的基于计数器的RCU}
\label{sec:app:toyrcu:Scalable Counter-Based RCU}

\Cref{lst:app:toyrcu:RCU Read-Side Using Per-Thread Reference-Count Pair}
(\path{rcu_rcpl.h})
展示了使用每线程引用计数器对的RCU实现的读端原语。
这个实现与
\cref{lst:app:toyrcu:RCU Read-Side Using Per-Thread Reference-Count}
中所示的非常相似，
主要区别在于添加了
\cref{lst:app:toyrcu:RCU Per-Thread Reference-Count Pair Data}
中所示的\co{rcu_refcnt}每线程数组。
与\cref{lst:app:toyrcu:RCU Read-Side Using Global Reference-Count Pair}
中的算法一样，
使用这个双元素数组可以防止读者使即使最不幸的更新者也饥饿。

\begin{listing}
\input{CodeSamples/defer/rcu_rcpl@define.fcv}
\caption{RCU每线程引用计数对数据}
\label{lst:app:toyrcu:RCU Per-Thread Reference-Count Pair Data}
\end{listing}

\begin{listing}
\input{CodeSamples/defer/rcu_rcpl@r.fcv}
\caption{使用每线程引用计数对的RCU读端}
\label{lst:app:toyrcu:RCU Read-Side Using Per-Thread Reference-Count Pair}
\end{listing}

\begin{listing}
\input{CodeSamples/defer/rcu_rcpl@u.fcv}
\caption{使用每线程引用计数对的RCU更新}
\label{lst:app:toyrcu:RCU Update Using Per-Thread Reference-Count Pair}
\end{listing}

\Cref{lst:app:toyrcu:RCU Update Using Per-Thread Reference-Count Pair}
(\path{rcu_rcpl.c})
展示了\co{synchronize_rcu()}的实现，以及一个名为\co{flip_counter_and_wait()}的辅助函数。
\begin{fcvref}[ln:defer:rcu_rcpl:u:sync]
\co{synchronize_rcu()}函数类似于
\cref{lst:app:toyrcu:RCU Update Using Global Reference-Count Pair}中所示的，
不同之处在于重复的计数器翻转被\clnref{flip1,flip2}上
对新辅助函数的一对调用所替代。
\end{fcvref}

\begin{fcvref}[ln:defer:rcu_rcpl:u:flip]
新的\co{flip_counter_and_wait()}函数在\clnref{atmset}上更新\co{rcu_idx}变量，
在\clnref{mb1}上执行内存屏障，
然后\clnrefrange{loop:b}{loop:e}
对每个线程的先前\co{rcu_refcnt}元素自旋，等待其变为零。
一旦所有这些元素都变为零，
它在\clnref{mb2}上执行另一个内存屏障并返回。
\end{fcvref}

这个实现仍然有几个缺点。
首先，需要翻转\co{rcu_idx}两次给更新带来了相当大的开销，
特别是当线程数量很多时。

其次，与以前一样，\co{synchronize_rcu()}现在必须检查数量随线程数
线性增长的变量，这给具有大量线程的应用程序带来了相当大的开销。

第三，与以前一样，尽管并发的RCU更新原则上可以由一个共同的宽限期来满足，
但这个实现将宽限期串行化了，阻止了宽限期共享。

第四，需要查阅\co{rcu_idx}使得读者比
\cref{sec:app:toyrcu:Scalable Single-Counter RCU}
中提出的实现稍慢。

最后，如文中所述，对每线程变量和枚举线程的需求在某些软件环境中可能有问题。

话虽如此，读端原语的扩展性非常好，
无论运行在单CPU还是64-CPU的\Power{5}系统上都只需要约115纳秒。
如上所述，\co{synchronize_rcu()}原语不能扩展，
开销从单个\Power{5} CPU上的接近一微秒到64-CPU系统上的接近200微秒不等。
可以想象这个实现可以作为生产级用户级RCU实现的基础。

下一节描述了一种允许更高效并发RCU更新的算法。

\section{支持共享宽限期的可扩展基于计数器的RCU}
\label{sec:app:toyrcu:Scalable Counter-Based RCU With Shared Grace Periods}

\Cref{lst:app:toyrcu:RCU Read-Side Using Per-Thread Reference-Count Pair and Shared Update}
(\path{rcu_rcpls.h})
展示了使用每线程引用计数对的RCU实现的读端原语，
与之前一样，但允许更新共享宽限期。
\begin{fcvref}[ln:defer:rcu_rcpls:r]
与\cref{lst:app:toyrcu:RCU Read-Side Using Per-Thread Reference-Count Pair}
中所示的早期实现的主要区别在于\co{rcu_idx}现在是一个自由计数的\co{long}，
因此\cref{lst:app:toyrcu:RCU Read-Side Using Per-Thread Reference-Count Pair and Shared Update}
的\clnref{lock:idx}必须屏蔽低位比特。
我们还从使用\co{atomic_read()}和\co{atomic_set()}
切换到使用\co{READ_ONCE()}。
数据也非常相似，如
\cref{lst:app:toyrcu:RCU Read-Side Using Per-Thread Reference-Count Pair and Shared Update Data}
所示，\co{rcu_idx}现在是\co{long}而不是\co{atomic_t}。
\end{fcvref}

\begin{listing}
\input{CodeSamples/defer/rcu_rcpls@define.fcv}
\caption{使用每线程引用计数对和共享更新数据的RCU读端数据}
\label{lst:app:toyrcu:RCU Read-Side Using Per-Thread Reference-Count Pair and Shared Update Data}
\end{listing}

\begin{listing}
\input{CodeSamples/defer/rcu_rcpls@r.fcv}
\caption{使用每线程引用计数对和共享更新的RCU读端}
\label{lst:app:toyrcu:RCU Read-Side Using Per-Thread Reference-Count Pair and Shared Update}
\end{listing}

\Cref{lst:app:toyrcu:RCU Shared Update Using Per-Thread Reference-Count Pair}
(\path{rcu_rcpls.c})
展示了\co{synchronize_rcu()}及其辅助函数\co{flip_counter_and_wait()}的实现。
这些与\cref{lst:app:toyrcu:RCU Update Using Per-Thread Reference-Count Pair}
中的类似。
\co{flip_counter_and_wait()}中的区别包括：
\begin{fcvref}[ln:defer:rcu_rcpls:u:flip]
\begin{enumerate}
\item	\Clnref{inc}使用\co{WRITE_ONCE()}而不是\co{atomic_set()}，
	并且是递增而不是取补。
\item	新的\clnref{mask}将计数器屏蔽到其最低位。
\end{enumerate}
\end{fcvref}

\begin{listing}
\input{CodeSamples/defer/rcu_rcpls@u.fcv}
\caption{使用每线程引用计数对的RCU共享更新}
\label{lst:app:toyrcu:RCU Shared Update Using Per-Thread Reference-Count Pair}
\end{listing}

\begin{fcvref}[ln:defer:rcu_rcpls:u:sync]
\co{synchronize_rcu()}的更改更为广泛：
\begin{enumerate}
\item	有一个新的\co{oldctr}局部变量，在\clnref{oldctr}上
	捕获获取锁之前的\co{rcu_idx}值。
\item	\Clnref{idx}使用\co{READ_ONCE()}而不是\co{atomic_read()}。
\item	\Clnrefrange{if:b}{ret}检查在获取锁期间是否有至少三次计数器翻转
	由其他线程执行，如果是，则释放锁，执行内存屏障并返回。
	在这种情况下，已经有两次完整的等待计数器归零，
	所以那些其他线程已经完成了所有必需的工作。
\item	在\clnrefrange{ifpair}{flip2}，只有当获取锁期间
	计数器翻转少于两次时，\co{flip_counter_and_wait()}才会被第二次调用。
	另一方面，如果有两次计数器翻转，某个其他线程已经做了
	一次完整的等待所有计数器归零，所以只需要再等一次。
\end{enumerate}
\end{fcvref}

使用这种方法，如果任意大量的线程同时调用\co{synchronize_rcu()}，
每个线程一个CPU，总共只会有三次等待计数器归零。

尽管有所改进，这个RCU实现仍然有一些缺点。
首先，与以前一样，需要翻转\co{rcu_idx}两次给更新带来了相当大的开销，
特别是当线程数量很多时。

其次，每个更新者仍然获取\co{rcu_gp_lock}，即使没有工作要做。
如果有大量并发更新，这可能导致严重的可扩展性限制。
有一些方法可以避免这个问题，
如在Linux内核的生产级实时RCU实现中所做的那样~\cite{PaulEMcKenney2007PreemptibleRCU}。

第三，这个实现需要每线程变量和枚举线程的能力，
这在某些软件环境中同样可能有问题。

最后，在32位机器上，给定的更新线程可能被抢占足够长的时间，
导致\co{rcu_idx}计数器溢出。
这可能导致该线程强制进行不必要的一对计数器翻转。
然而，即使每个宽限期只需要一微秒，
违规线程也需要被抢占超过一个小时，
在这种情况下，额外的一对计数器翻转可能是你最不需要担心的问题。

与\cref{sec:app:toyrcu:Simple Counter-Based RCU}中描述的实现一样，
读端原语的扩展性极好，
无论CPU数量如何都只有约115纳秒的开销。
\co{synchronize_rcu()}原语仍然很昂贵，
从约一微秒到约16微秒不等。
尽管如此，这比\cref{sec:app:toyrcu:Scalable Counter-Based RCU}
中实现所产生的约200微秒要便宜得多。
因此，尽管有这些缺点，可以想象这个RCU实现在实际应用中用于生产。

\QuickQuiz{
	所有这些玩具RCU实现要么在\co{rcu_read_lock()}和\co{rcu_read_unlock()}中
	有原子操作，要么\co{synchronize_rcu()}的开销随线程数线性增长。
	在什么情况下，RCU实现可以对所有这三个原语都有轻量级实现，
	且都具有确定性（$\O{1}$）的开销和延迟？
}\QuickQuizAnswer{
	专用的单处理器RCU实现可以达到这个理想状态~\cite{PaulEMcKenney2009BloatwatchRCU}。
}\QuickQuizEnd

回顾
\cref{lst:app:toyrcu:RCU Read-Side Using Per-Thread Reference-Count Pair and Shared Update}，
我们看到有一次全局变量访问和不少于四次线程局部变量访问。
鉴于在实现POSIX线程的系统上线程局部访问的相对高成本，
很容易想到将三个线程局部变量合并为一个结构，
允许\co{rcu_read_lock()}和\co{rcu_read_unlock()}
通过一次线程局部存储访问来访问其线程局部数据。
然而，更好的方法是将线程局部访问的次数减少到一次，
这将在下一节中完成。

\section{基于自由运行计数器的RCU}
\label{sec:app:toyrcu:RCU Based on Free-Running Counter}

\Cref{lst:app:toyrcu:Free-Running Counter Using RCU}
(\path{rcu.h}和\path{rcu.c})
展示了一个基于单一全局自由运行计数器的RCU实现，
该计数器只取偶数值，数据如
\cref{lst:app:toyrcu:Data for Free-Running Counter Using RCU}所示。

\begin{listing}
\input{CodeSamples/defer/rcu@define.fcv}
\caption{使用自由运行计数器的RCU数据}
\label{lst:app:toyrcu:Data for Free-Running Counter Using RCU}
\end{listing}

\begin{listing}
\input{CodeSamples/defer/rcu@read_lock_unlock.fcv}\vspace*{-11pt}\fvset{firstnumber=last}
\input{CodeSamples/defer/rcu@synchronize.fcv}\fvset{firstnumber=auto}
\caption{使用自由运行计数器的RCU}
\label{lst:app:toyrcu:Free-Running Counter Using RCU}
\end{listing}

生成的\co{rcu_read_lock()}实现极其简单直接。
\begin{fcvref}[ln:defer:rcu:read_lock_unlock:lock]
\Clnref{gp1,gp2}简单地将值一加到全局自由运行的\co{rcu_gp_ctr}变量上，
并将生成的奇数值存储到\co{rcu_reader_gp}每线程变量中。
\Clnref{mb}执行一个\IX{memory barrier}以防止后续RCU读端临界区的内容"泄漏出去"。
\end{fcvref}

\begin{fcvref}[ln:defer:rcu:read_lock_unlock:unlock]
\co{rcu_read_unlock()}的实现类似。
\Clnref{mb}执行一个内存屏障，同样是为了防止先前的RCU读端临界区"泄漏出去"。
\Clnref{gp1,gp2}然后将\co{rcu_gp_ctr}全局变量复制到
\co{rcu_reader_gp}每线程变量中，使这个每线程变量留下一个偶数值，
这样\co{synchronize_rcu()}的并发实例就知道要忽略它。
\end{fcvref}

\QuickQuiz{
	\begin{fcvref}[ln:defer:rcu:read_lock_unlock:unlock]
	如果任何偶数值都足以告诉\co{synchronize_rcu()}忽略给定的任务，
	为什么\cref{lst:app:toyrcu:Free-Running Counter Using RCU}的
	\clnref{gp1,gp2}不简单地将零赋给\co{rcu_reader_gp}？
	\end{fcvref}
}\QuickQuizAnswer{
	赋值零（或任何其他偶数常量）确实可以工作，
	但赋值\co{rcu_gp_ctr}的值可以提供有价值的调试辅助，
	因为它给开发者一个关于对应线程上次退出RCU读端临界区时间的概念。
}\QuickQuizEnd

\begin{fcvref}[ln:defer:rcu:synchronize:syn]
因此，\co{synchronize_rcu()}可以等待所有每线程\co{rcu_reader_gp}变量
取偶数值。
然而，可以做得比这好得多，因为\co{synchronize_rcu()}只需要等待
\emph{先前存在的}RCU读端临界区。
\Clnref{mb1}执行一个内存屏障，以防止先前对RCU保护的数据结构的操作
被重排序（无论是由CPU还是编译器）到\clnref{increasegp}上的递增之后。
\Clnref{spinlock}获取\co{rcu_gp_lock}
（\clnref{spinunlock}释放它），
以防止多个\co{synchronize_rcu()}实例并发运行。
\Clnref{increasegp}然后将全局\co{rcu_gp_ctr}变量递增二，
使得所有先前存在的RCU读端临界区将具有对应的每线程\co{rcu_reader_gp}变量，
其值小于\co{rcu_gp_ctr}的值（以机器字长为模）。
还请回忆，\co{rcu_reader_gp}值为偶数的线程不在RCU读端临界区中，
因此\clnrefrange{scan:b}{scan:e}
扫描\co{rcu_reader_gp}值，直到它们全部要么是偶数（\clnref{even}），
要么大于全局\co{rcu_gp_ctr}（\clnrefrange{gt1}{gt2}）。
\Clnref{poll}阻塞一小段时间等待先前存在的RCU读端临界区，
但如果\IXh{grace-period}{latency}至关重要，可以用自旋循环替代。
最后，\clnref{mb3}上的内存屏障确保任何后续的销毁不会被重排序到前面的循环中。
\end{fcvref}

\QuickQuiz{
	\begin{fcvref}[ln:defer:rcu:synchronize:syn]
	为什么\cref{lst:app:toyrcu:Free-Running Counter Using RCU}
	的\clnref{mb1,mb3}上需要内存屏障？
	\clnref{spinlock,spinunlock}上锁原语中固有的内存屏障不够吗？
	\end{fcvref}
}\QuickQuizAnswer{
	这些内存屏障是必需的，因为锁原语只保证限定临界区。
	锁原语完全没有义务防止其他代码泄漏到临界区中。
	因此，这对内存屏障是必需的，以防止这种代码移动，
	无论是由编译器还是由CPU执行的\@。
}\QuickQuizEnd

这种方法实现了更好的读端性能，
无论\Power{5} CPU的数量如何，都只产生约63纳秒的开销。
更新产生更多开销，从单个\Power{5} CPU上的约500纳秒
到64个这样的CPU上的超过100\emph{微秒}不等。

\QuickQuiz{
	\cref{sec:app:toyrcu:Scalable Counter-Based RCU With Shared Grace Periods}
	中描述的更新端批处理优化能否应用于
	\cref{lst:app:toyrcu:Free-Running Counter Using RCU}中所示的实现？
}\QuickQuizAnswer{
	确实可以，只需做一些修改。
	这项工作留给读者作为练习。
}\QuickQuizEnd

\begin{fcvref}[ln:defer:rcu:read_lock_unlock:lock]
除了前面提到的高更新端开销外，这个实现还有一些严重的缺点。
首先，不再允许嵌套RCU读端临界区，这个话题将在下一节中讨论。
其次，如果读者在\cref{lst:app:toyrcu:Free-Running Counter Using RCU}
的\clnref{gp1}处被抢占，
已经从\co{rcu_gp_ctr}取值但尚未存储到\co{rcu_reader_gp}，
而\co{rcu_gp_ctr}计数器随后经过了其可能值的一半以上但不到全部，
那么\co{synchronize_rcu()}将忽略后续的RCU读端临界区。
第三也是最后，这个实现要求封装的软件环境能够枚举线程并维护每线程变量。
\end{fcvref}

\QuickQuiz{
	\begin{fcvref}[ln:defer:rcu:read_lock_unlock:lock]
	读者在\cref{lst:app:toyrcu:Free-Running Counter Using RCU}
	的\clnrefrange{gp1}{gp2}处被抢占的可能性
	是一个真正的问题吗？换句话说，是否存在一个真正的事件序列
	可能导致失败？
	如果不是，为什么不是？
	如果是，事件序列是什么，如何解决这个失败？
	\end{fcvref}
}\QuickQuizAnswer{
	这是一个真正的问题，存在导致失败的事件序列，
	并且有多种可能的解决方法。
	有关更多细节，请参阅
	\cref{sec:app:toyrcu:Nestable RCU Based on Free-Running Counter}
	末尾附近的小测验。
	将讨论放在那里的原因是(1)给你更多时间思考，
	(2)因为该节中添加的嵌套支持大大减少了计数器溢出所需的时间。
}\QuickQuizEnd

\section{基于自由运行计数器的可嵌套RCU}
\label{sec:app:toyrcu:Nestable RCU Based on Free-Running Counter}

\Cref{lst:app:toyrcu:Nestable RCU Using a Free-Running Counter}
(\path{rcu_nest.h}和\path{rcu_nest.c})
展示了一个基于单一全局自由运行计数器的RCU实现，
但允许RCU读端临界区的嵌套。
这种可嵌套性通过保留全局\co{rcu_gp_ctr}的低位来计数嵌套来实现，
使用\cref{lst:app:toyrcu:Data for Nestable RCU Using a Free-Running Counter}
中所示的定义。
这是\cref{sec:app:toyrcu:RCU Based on Free-Running Counter}
中方案的推广，后者可以被认为保留了单个低位来计数嵌套深度。
两个C预处理器宏用于实现这一安排，
\co{RCU_GP_CTR_NEST_MASK}和\co{RCU_GP_CTR_BOTTOM_BIT}。
它们之间的关系是：\co{RCU_GP_CTR_NEST_MASK=RCU_GP_CTR_BOTTOM_BIT-1}。
\co{RCU_GP_CTR_BOTTOM_BIT}宏包含一个位，位于为计数嵌套而保留的位的正上方，
\co{RCU_GP_CTR_NEST_MASK}的所有一位覆盖了\co{rcu_gp_ctr}中
用于计数嵌套的区域。
显然，这两个C预处理器宏必须保留计数器的足够多低位，
以允许RCU读端临界区的最大所需嵌套，
这个实现保留了七位，最大RCU读端临界区嵌套深度为127，
这应该远远超出大多数应用程序所需的。

\begin{listing}
\input{CodeSamples/defer/rcu_nest@define.fcv}
\caption{使用自由运行计数器的可嵌套RCU数据}
\label{lst:app:toyrcu:Data for Nestable RCU Using a Free-Running Counter}
\end{listing}

\begin{listing}
\input{CodeSamples/defer/rcu_nest@read_lock_unlock.fcv}\vspace*{-11pt}\fvset{firstnumber=last}
\input{CodeSamples/defer/rcu_nest@synchronize.fcv}\fvset{firstnumber=auto}
\caption{使用自由运行计数器的可嵌套RCU}
\label{lst:app:toyrcu:Nestable RCU Using a Free-Running Counter}
\end{listing}

\begin{fcvref}[ln:defer:rcu_nest:read_lock_unlock:lock]
生成的\co{rcu_read_lock()}实现仍然相当简单直接。
\Clnref{readgp}将指向该线程\co{rcu_reader_gp}实例的指针
放入局部变量\co{rrgp}中，最大限度地减少了对pthreads线程局部状态API的
昂贵调用次数\@。
\Clnref{wtmp1}将\co{rcu_reader_gp}的当前值记录到另一个局部变量\co{tmp}中，
\clnref{checktmp}检查低位是否为零，如果是则表明这是最外层的\co{rcu_read_lock()}。
如果是这样，\clnref{wtmp2}将全局\co{rcu_gp_ctr}放入\co{tmp}中，
因为\clnref{wtmp1}先前获取的当前值可能已经过时。
无论哪种情况，\clnref{inctmp}递增嵌套深度，
你会记得它存储在计数器的七个低位中。
\Clnref{writegp}将更新后的计数器存回该线程的\co{rcu_reader_gp}实例中，
最后，\clnref{mb1}执行内存屏障，
以防止RCU读端临界区泄漏到调用\co{rcu_read_lock()}之前的代码中。
\end{fcvref}

换句话说，这个\co{rcu_read_lock()}实现获取全局\co{rcu_gp_ctr}的副本，
除非当前的\co{rcu_read_lock()}调用嵌套在RCU读端临界区内，
在这种情况下，它转而获取当前线程\co{rcu_reader_gp}实例的内容。
无论哪种方式，它递增所获取的任何值以记录额外的嵌套级别，
并将结果存储在当前线程的\co{rcu_reader_gp}实例中。

\begin{fcvref}[ln:defer:rcu_nest:read_lock_unlock:unlock]
有趣的是，尽管\co{rcu_read_lock()}有差异，
\co{rcu_read_unlock()}的实现与
\cref{sec:app:toyrcu:RCU Based on Free-Running Counter}
中所示的大体相似。
\Clnref{mb1}执行内存屏障，以防止RCU读端临界区
泄漏到调用\co{rcu_read_unlock()}之后的代码中，
\clnref{decgp}递减该线程的\co{rcu_reader_gp}实例，
其效果是递减\co{rcu_reader_gp}低位中包含的嵌套计数。
这个原语的调试版本会检查（在递减之前！\@）这些低位是否非零。
\end{fcvref}

\begin{fcvref}[ln:defer:rcu_nest:synchronize:syn]
\co{synchronize_rcu()}的实现与
\cref{sec:app:toyrcu:RCU Based on Free-Running Counter}
中所示的非常相似。
有两个区别。
第一个是\clnref{incgp1,incgp2}将\co{RCU_GP_CTR_BOTTOM_BIT}
加到全局\co{rcu_gp_ctr}上，而不是加常量"2"，
第二个是\clnref{ongoing}上的比较被抽象到一个单独的函数中，
该函数检查\co{RCU_GP_CTR_NEST_MASK}指示的位，
而不是无条件检查低位。
\end{fcvref}

这种方法实现了几乎等于
\cref{sec:app:toyrcu:RCU Based on Free-Running Counter}
中所示的读端性能，无论\Power{5} CPU的数量如何，都只产生约65纳秒的开销。
更新再次产生更多开销，从单个\Power{5} CPU上的约600纳秒
到64个这样的CPU上的超过100\emph{微秒}不等。

\QuickQuiz{
	为什么不像前一节那样简单地维护一个单独的每线程嵌套级别变量，
	而不是进行所有这些复杂的位操作？
}\QuickQuizAnswer{
	单独的每线程变量的表面简单性是一个误导。
	这种方法在仔细的操作排序方面会带来更大的复杂性，
	特别是如果要允许信号处理程序包含RCU读端临界区的话。
	但不要只听我说，自己编码试试看最终会得到什么！
}\QuickQuizEnd

这个实现与
\cref{sec:app:toyrcu:RCU Based on Free-Running Counter}
中的具有相同的缺点，只是现在允许嵌套RCU读端临界区。
此外，在32位系统上，这种方法缩短了全局\co{rcu_gp_ctr}变量溢出所需的时间。
下一节展示了一种大大增加溢出发生所需时间的方法，
同时大大减少读端开销。

\QuickQuizSeries{%
\QuickQuizB{
	给定\cref{lst:app:toyrcu:Nestable RCU Using a Free-Running Counter}
	中所示的算法，你如何将全局\co{rcu_gp_ctr}溢出所需的时间翻倍？
}\QuickQuizAnswerB{
	\begin{fcvref}[ln:defer:rcu_nest:synchronize:syn]
	一种方法是将\clnref{lt1,lt2}上的大小比较替换为
	每线程\co{rcu_reader_gp}变量与
	\co{rcu_gp_ctr+RCU_GP_CTR_BOTTOM_BIT}的不等式检查。
	\end{fcvref}
}\QuickQuizEndB
%
\QuickQuizE{
	再次，给定\cref{lst:app:toyrcu:Nestable RCU Using a Free-Running Counter}
	中所示的算法，计数器溢出是否是致命的？
	为什么是或为什么不是？
	如果是致命的，可以做什么来修复？
}\QuickQuizAnswerE{
	它确实可能是致命的。
	要理解这一点，考虑以下事件序列：
	\begin{enumerate}
	\item	线程0进入\co{rcu_read_lock()}，确定它不是嵌套的，
		因此获取全局\co{rcu_gp_ctr}的值。
		然后线程0被抢占了极长的时间
		（在存储到其每线程\co{rcu_reader_gp}变量之前）。
	\item	其他线程反复调用\co{synchronize_rcu()}，
		使得全局\co{rcu_gp_ctr}的新值现在比线程0获取时
		少了\co{RCU_GP_CTR_BOTTOM_BIT}。
	\item	线程0现在开始重新运行，并存储到其每线程\co{rcu_reader_gp}变量中。
		它存储的值比全局\co{rcu_gp_ctr}大\co{RCU_GP_CTR_BOTTOM_BIT+1}。
	\item	线程0获取了对RCU保护的数据元素A的引用。
	\item	线程1现在移除了线程0刚刚获取引用的数据元素A。
	\item	线程1调用\co{synchronize_rcu()}，
		将全局\co{rcu_gp_ctr}递增\co{RCU_GP_CTR_BOTTOM_BIT}。
		然后它检查所有每线程\co{rcu_reader_gp}变量，
		但线程0的值（错误地）表明它是在线程1调用\co{synchronize_rcu()}
		之后开始的，所以线程1不等待线程0完成其RCU读端临界区。
	\item	线程1然后释放数据元素A，而线程0仍在引用它。
	\end{enumerate}

	请注意，这种情况也可能发生在
	\cref{sec:app:toyrcu:RCU Based on Free-Running Counter}
	中提出的实现中。

	修复这个问题的一种策略是使用64位计数器，
	使得溢出所需的时间超过计算机系统的有用寿命。
	请注意，32位x86 CPU系列的非古董成员允许通过
	\co{cmpxchg64b}指令原子操作64位计数器。

	另一种策略是限制允许宽限期发生的速率，以达到类似的效果。
	例如，\co{synchronize_rcu()}可以记录上次被调用的时间，
	任何后续调用都会检查这个时间并根据需要阻塞以强制所需的间隔。
	例如，如果计数器的低四位被保留用于嵌套，
	并且宽限期被允许每秒最多发生十次，
	那么计数器溢出需要超过300天。
	然而，如果存在任何可能性使系统在整整300天内
	被CPU密集型高优先级实时线程完全占满，这种方法就没有帮助。
	（也许是一个很小的可能性，但最好提前考虑。）

	第三种方法是从管理上禁止相关系统中的实时线程。
	在这种情况下，被抢占的进程会随着时间提升优先级，
	因此在计数器有机会溢出之前就能运行。
	当然，对于实时应用程序，这种方法帮助不大。

	第四种方法是让\co{rcu_read_lock()}在存储到其每线程\co{rcu_reader_gp}
	计数器之后重新检查全局\co{rcu_gp_ctr}的值，
	如果全局\co{rcu_gp_ctr}的新值不合适则重试。
	这可以工作，但在\co{rcu_read_lock()}中引入了非确定性执行时间。
	另一方面，如果你的应用程序被抢占的时间长到足以使计数器溢出，
	你无论如何都没有确定性执行时间的希望！

	第五种方法是让宽限期进程等待所有读者意识到新的宽限期。
	这在理论上运行得很好，但如果读者在RCU读端临界区外无限期阻塞则会挂起。

	最后一种方法是，奇怪的是，使用单比特宽限期计数器，
	并让每次\co{synchronize_rcu()}调用对其算法进行两遍。
	这是用户空间RCU~\cite{MathieuDesnoyers2009URCU}所使用的方法，
	在期刊文章和补充材料~\cite[Appendix D]{MathieuDesnoyers2012URCU}
	中有详细描述。
}\QuickQuizEndE
}

\section{基于静默状态的RCU}
\label{sec:app:toyrcu:RCU Based on Quiescent States}

\begin{fcvref}[ln:defer:rcu_qs:read_lock_unlock]
\Cref{lst:app:toyrcu:Quiescent-State-Based RCU Read Side}
(\path{rcu_qs.h})
展示了用于构建基于\IXpl{quiescent state}的用户级RCU实现的读端原语，
数据如\cref{lst:app:toyrcu:Data for Quiescent-State-Based RCU}所示。
从清单中的\clnrefrange{lock:b}{unlock:e}可以看出，
\co{rcu_read_lock()}和\co{rcu_read_unlock()}原语什么都不做，
事实上可以预期它们会被内联并优化掉，就像在Linux内核的服务器构建版本中一样。
这是因为基于静默状态的RCU实现
使用上述静默状态来\emph{近似}RCU读端临界区的范围。
每个静默状态包含对\co{rcu_quiescent_state()}的调用，
如清单中\clnrefrange{qs:b}{qs:e}所示。
进入扩展静默状态的线程（例如，阻塞时）
可以在进入扩展静默状态时调用\co{rcu_thread_offline()}
（\clnrefrange{offline:b}{offline:e}），
然后在离开时调用\co{rcu_thread_online()}
（\clnrefrange{online:b}{online:e}）。
因此，\co{rcu_thread_online()}类似于\co{rcu_read_lock()}，
\co{rcu_thread_offline()}类似于\co{rcu_read_unlock()}。
此外，\co{rcu_quiescent_state()}可以被认为是
\co{rcu_thread_online()}紧接着\co{rcu_thread_offline()}。\footnote{
	虽然清单中的代码与\co{rcu_quiescent_state()}等同于
	\co{rcu_thread_online()}紧接着\co{rcu_thread_offline()}一致，
	但这种关系被性能优化所掩盖。}
从RCU读端临界区内调用\co{rcu_quiescent_state()}、
\co{rcu_thread_offline()}或\co{rcu_thread_online()}是非法的。
\end{fcvref}

\begin{listing}
\input{CodeSamples/defer/rcu_qs@define.fcv}
\caption{基于静默状态的RCU数据}
\label{lst:app:toyrcu:Data for Quiescent-State-Based RCU}
\end{listing}

\begin{listing}
\input{CodeSamples/defer/rcu_qs@read_lock_unlock.fcv}
\caption{基于静默状态的RCU读端}
\label{lst:app:toyrcu:Quiescent-State-Based RCU Read Side}
\end{listing}

\begin{fcvref}[ln:defer:rcu_qs:read_lock_unlock:qs]
在\co{rcu_quiescent_state()}中，\clnref{mb1}执行内存屏障，
以防止静默状态之前的任何代码（包括可能的RCU读端临界区）
被重排序到静默状态中。
\Clnrefrange{gp1}{gp2}获取全局\co{rcu_gp_ctr}的副本，
使用\co{READ_ONCE()}确保编译器不会采用任何导致\co{rcu_gp_ctr}
被获取超过一次的优化，
然后将获取的值加一并存储到每线程\co{rcu_reader_qs_gp}变量中，
这样\co{synchronize_rcu()}的任何并发实例将看到奇数值，
从而意识到新的RCU读端临界区已经开始。
等待旧的RCU读端临界区的\co{synchronize_rcu()}实例
因此会知道忽略这个新的。
最后，\clnref{mb2}执行内存屏障，防止后续代码
（包括可能的RCU读端临界区）与\clnrefrange{gp1}{gp2}重排序。
\end{fcvref}

\QuickQuiz{
	\begin{fcvref}[ln:defer:rcu_qs:read_lock_unlock:qs]
	\cref{lst:app:toyrcu:Quiescent-State-Based RCU Read Side}
	的\clnref{mb2}上显示的额外内存屏障
	不是大大增加了\co{rcu_quiescent_state()}的开销吗？
	\end{fcvref}
}\QuickQuizAnswer{
	\begin{fcvref}[ln:defer:rcu_qs:read_lock_unlock:qs]
	确实如此！
	因此，使用这个RCU实现的应用程序应该节制地调用\co{rcu_quiescent_state()}，
	而大多数时候使用\co{rcu_read_lock()}和\co{rcu_read_unlock()}。

	然而，这个内存屏障是绝对必需的，
	以确保其他线程在调用者执行的任何后续RCU读端临界区之前
	看到\clnrefrange{gp1}{gp2}上的存储。
	\end{fcvref}
}\QuickQuizEnd

某些应用程序可能只是偶尔使用RCU，但使用时会非常频繁。
这类应用程序可能选择在开始使用RCU时调用\co{rcu_thread_online()}，
在不再使用RCU时调用\co{rcu_thread_offline()}\@。
从\co{rcu_thread_offline()}调用到后续\co{rcu_thread_online()}调用
之间的时间是一个扩展静默状态，
因此RCU不会期望在此期间注册显式的静默状态。

\co{rcu_thread_offline()}函数简单地将每线程\co{rcu_reader_qs_gp}变量
设置为\co{rcu_gp_ctr}的当前值，该值是偶数。
\co{synchronize_rcu()}的任何并发实例因此将知道忽略该线程。

\QuickQuiz{
	\begin{fcvref}[ln:defer:rcu_qs:read_lock_unlock:qs]
	为什么\cref{lst:app:toyrcu:Quiescent-State-Based RCU Read Side}
	的\clnref{mb1,mb2}上需要两个内存屏障？
	\end{fcvref}
}\QuickQuizAnswer{
	\begin{fcvref}[ln:defer:rcu_qs:read_lock_unlock:qs]
	\clnref{mb1}上的内存屏障防止\co{rcu_thread_offline()}调用之前
	可能存在的任何RCU读端临界区被编译器或CPU重排序到
	\clnrefrange{gp1}{gp2}上的赋值之后。
	严格来说，\clnref{mb2}上的内存屏障是不必要的，
	因为在\co{rcu_thread_offline()}调用之后
	拥有任何RCU读端临界区是非法的。
	\end{fcvref}
}\QuickQuizEnd

\co{rcu_thread_online()}函数简单地调用\co{rcu_quiescent_state()}，
从而标记扩展静默状态的结束。

\begin{listing}
\input{CodeSamples/defer/rcu_qs@synchronize.fcv}
\caption{使用静默状态的RCU更新端}
\label{lst:app:toyrcu:RCU Update Side Using Quiescent States}
\end{listing}

\Cref{lst:app:toyrcu:RCU Update Side Using Quiescent States}
(\path{rcu_qs.c})
展示了\co{synchronize_rcu()}的实现，它与前面各节的实现非常相似。

这个实现具有极快的读端原语，
\co{rcu_read_lock()}--\co{rcu_read_unlock()}往返只产生约50\emph{皮秒}的开销。
\co{synchronize_rcu()}的开销从单CPU \Power{5}系统上的约600纳秒
到64-CPU系统上的超过100微秒不等。

\QuickQuiz{
	诚然，2008年\Power{}系统的时钟频率相当高，
	但即使是5\,GHz的时钟频率也不足以在50皮秒内执行循环！
	这是怎么回事？
}\QuickQuizAnswer{
	由于测量循环包含一对空函数，编译器将其优化掉了。
	测量循环在每次调用\co{rcu_quiescent_state()}之间运行1,000遍，
	因此这个测量值大约是单次调用\co{rcu_quiescent_state()}开销的千分之一。
}\QuickQuizEnd

然而，这个实现要求每个线程要么周期性地调用\co{rcu_quiescent_state()}，
要么为扩展静默状态调用\co{rcu_thread_offline()}。
需要周期性地调用这些函数使得这个实现在某些情况下难以使用，
例如某些类型的库函数。

\QuickQuizSeries{%
\QuickQuizB{
	为什么代码在库中这一事实会对使用
	\cref{lst:app:toyrcu:Quiescent-State-Based RCU Read Side,%
	lst:app:toyrcu:RCU Update Side Using Quiescent States}
	中所示RCU实现的难易程度产生影响？
}\QuickQuizAnswerB{
	库函数对调用者完全没有控制权，
	因此无法强制调用者周期性地调用\co{rcu_quiescent_state()}。
	另一方面，一个对给定RCU保护的数据结构进行大量引用的库函数
	可能能够在进入时调用\co{rcu_thread_online()}，
	周期性地调用\co{rcu_quiescent_state()}，
	并在退出时调用\co{rcu_thread_offline()}。
}\QuickQuizEndB
%
\QuickQuizE{
	但如果你在调用\co{synchronize_rcu()}期间持有一个锁，
	然后在RCU读端临界区内获取同一个锁呢？
	这应该是死锁，但一个完全不生成代码的原语怎么可能参与死锁环？
}\QuickQuizAnswerE{
	请注意，RCU读端临界区实际上被扩展到了包围的
	\co{rcu_read_lock()}和\co{rcu_read_unlock()}之外，
	延伸到前一个和下一个\co{rcu_quiescent_state()}调用。
	这个\co{rcu_quiescent_state()}可以被认为是
	\co{rcu_read_unlock()}紧接着\co{rcu_read_lock()}。

	即便如此，实际的死锁本身将涉及RCU读端临界区中的锁获取
	和\co{synchronize_rcu()}，而不会涉及\co{rcu_quiescent_state()}。
}\QuickQuizEndE
}

此外，这个实现不允许并发的\co{synchronize_rcu()}调用共享宽限期。
话虽如此，可以很容易地想象基于这个版本的RCU构建一个生产级RCU实现\@。

\section{玩具RCU实现总结}
\label{sec:app:toyrcu:Summary of Toy RCU Implementations}

如果你读到了这里，恭喜！
你现在应该对RCU本身以及其封装软件环境和应用程序的需求有了更清晰的理解。
希望获得更深入理解的人可以阅读生产级RCU实现的
描述~\cite{MathieuDesnoyers2012URCU,PaulEMcKenney2007PreemptibleRCU,PaulEMcKenney2008HierarchicalRCU,PaulEMcKenney2009BloatwatchRCU}。

前面各节列出了各种RCU原语的一些期望属性。
以下列表供希望创建新RCU实现的人方便参考。

\begin{enumerate}
\item	必须有读端原语（如\co{rcu_read_lock()}和\co{rcu_read_unlock()}）
	和宽限期原语（如\co{synchronize_rcu()}和\co{call_rcu()}），
	使得在宽限期开始时存在的任何RCU读端临界区
	在宽限期结束时已经完成。
\item	RCU读端原语应该有最小的开销。
	特别是，应避免昂贵的操作，如缓存未命中、
	原子指令、内存屏障和分支。
\item	RCU读端原语应该具有$\O{1}$的计算复杂度以支持实时使用。
	这意味着读者与更新者并发运行。
	（但请注意，可抢占的Linux内核RCU的\co{rcu_read_unlock()}
	对于低优先级线程和极长运行的RCU读端临界区结束时
	可能违反这一期望属性。）
\item	RCU读端原语应该在所有上下文中可用
	（在Linux内核中，它们在除了空闲循环深处、离线CPU和
	内核进入/退出代码路径之外的所有地方都被允许）。
	一个重要的特殊情况是RCU读端原语可在RCU读端临界区内使用，
	换句话说，可以嵌套RCU读端临界区。
\item	RCU读端原语应该是无条件的，没有失败返回。
	这个属性极其重要，因为失败检查增加了复杂性并使测试和验证更加复杂。
\item	除了静默状态（因此也是宽限期）之外的任何操作
	都应该在RCU读端临界区中被允许。
	特别是，不可撤销的操作如I/O应该被允许。
\item	应该可以在RCU读端临界区内执行时更新RCU保护的数据结构。
\item	RCU读端和更新端原语都应该独立于内存分配器的设计和实现，
	换句话说，同一个RCU实现应该能够保护给定的数据结构，
	而不管数据元素是如何分配和释放的。
\item	RCU宽限期不应该被在RCU读端临界区外停止的线程阻塞。
	（但请注意，大多数基于静默状态的实现违反了这一期望属性。）
\end{enumerate}

\QuickQuiz{
	既然宽限期在RCU读端临界区内是被禁止的，
	那么在RCU读端临界区内如何才能更新RCU数据结构？
}\QuickQuizAnswer{
	这种情况是异步宽限期原语（如\co{call_rcu()}）存在的原因之一。
	这个原语可以在RCU读端临界区内调用，
	指定的RCU回调将在宽限期经过后的稍后时间被调用。

	在RCU读端临界区内执行RCU更新的能力可以非常方便，
	类似于读写锁的（神话般的）无条件读到写升级。
}\QuickQuizEnd

希望了解更多Linux内核RCU需求细节的人应参考其内部
文档~\cite{PaulEMcKenney2015rcu-requirements}。

\QuickQuizAnswersChp{qqztoyrcu}
